\documentclass[11pt,a4paper,twocolumn]{article}

\usepackage[T2A]{fontenc}
\usepackage[utf8]{inputenc}
\usepackage[russian]{babel}
\usepackage[a4paper,top=2.2cm,bottom=2.2cm,left=1.8cm,right=1.8cm,columnsep=0.8cm,headheight=28pt]{geometry}
\usepackage{graphicx}
\usepackage{amsmath}
\usepackage{booktabs}
\usepackage{longtable}
\usepackage{array}
\usepackage{pdflscape}
\usepackage{fancyhdr}
\usepackage{titlesec}
\usepackage{xcolor}
\usepackage[style=authoryear,backend=biber,maxcitenames=2,mincitenames=1,maxbibnames=3,uniquelist=false,uniquename=false,dashed=true]{biblatex}
\usepackage[autostyle]{csquotes}
\usepackage[hidelinks]{hyperref}
\usepackage{url}

\addbibresource{references.bib}
\renewcommand*{\bibfont}{\small}
\DefineBibliographyStrings{russian}{andothers={и~др\adddot}}

\newcommand{\var}[1]{\texttt{#1}}
\newcommand{\head}[1]{\textbf{#1}}

\titleformat{\section}{\Large\bfseries}{\thesection}{1em}{}
\titleformat{\subsection}{\large\bfseries}{\thesubsection}{1em}{}
\titlespacing*{\paragraph}{0pt}{0.8ex plus 0.3ex}{0.6em}

\pagestyle{fancy}
\fancyhf{}
\fancyhead[L]{\itshape Научная студия 2026. Вычислительная онкология. Общий анализ данных.}
\fancyfoot[C]{\thepage}
\renewcommand{\headrulewidth}{0.4pt}

\fancypagestyle{first}{%
  \fancyhf{}
  \fancyhead[L]{\raisebox{-0.35\height}{\includegraphics[height=26pt]{logo.png}}}
  \fancyhead[R]{\itshape\small Октябрь, 2026\\ В.~Алексеев, С.~Зинина, И.~Коновалов, В.~Горбатенко, А.~Соколова}
  \fancyfoot[C]{\thepage}
  \renewcommand{\headrulewidth}{0pt}}

\setlength{\parskip}{0pt}
\emergencystretch=1.5em

\begin{document}
\thispagestyle{first}

\setcounter{section}{1}
\section{Материалы и методы}
\setcounter{subsection}{2}

\subsection{Исход и момент применения модели}

\paragraph{Целевая переменная.}
Модель предсказывает локальный рецидив первичной опухоли после радикальной лучевой терапии (ЛТ). Целевая переменная \var{Target\_local\_recurrence} равна 1, если за период наблюдения зарегистрирован рецидив в той же или иной подобласти ротоглотки, и 0, если рецидива не было; значение перенесено без преобразований из переменной \var{Local\_tumor\_recurrence}, заполненной авторами набора по данным наблюдения \parencite{elhalawani2017}. Эта конечная точка выбрана потому, что локальный контроль~--- прямая цель облучения первичного очага, а радиомические признаки описывают именно его (GTVp).

\paragraph{Горизонт прогноза.}
В таблицах нет дат, длительности наблюдения и времени до рецидива, поэтому исход нельзя пересчитать на иной горизонт и он используется в форме, заданной организаторами. Горизонт задан критерием включения в когорту~--- наблюдением более двух лет \parencite{elhalawani2017}: этот срок покрывает период, на который приходится основная часть локальных неудач после ЛТ. По критериям отбора все 128 пациентов без события наблюдались не меньше горизонта, но индивидуально проверить это по таблице нельзя, что является ограничением работы.

\paragraph{Неопределённый исход.}
У 158 пациентов тестовой части исход скрыт организаторами (пустая ячейка), т.\,е. неизвестен, а не отрицателен. Присвоение им нуля занизило бы частоту рецидива и внесло ошибку разметки, поэтому они исключены целиком. В анализ вошли 140 пациентов обучающей части: 12 с рецидивом (8,6~\%) и 128 без него. Малое число событий~--- главное ограничение данных, требующее жёсткого ограничения числа признаков на этапе моделирования \parencite{riley2020}.

\paragraph{Момент применения модели.}
Модель применяется после стадирования и утверждения плана лечения, до начала любой его части. К этому моменту известны диагностическая КТ с контрастированием, по которой построена маска GTVp (по критериям отбора она выполнена до лечения), демографические данные, анамнез курения, TNM и стадия, степень дифференцировки по биопсии, предписанные доза и фракционирование и решение о химиотерапии. Оценка риска при уже выбранном режиме позволяет обсуждать эскалацию дозы или более плотное наблюдение после ЛТ; раньше параметры лечения неизвестны, позже прогноз теряет практическую ценность. Поэтому как признаки нельзя использовать фактическую длительность курса ЛТ \var{Radiation\_treatment\_course\_duration}, зависящую от перерывов и известную лишь после окончания лечения, и жизненный статус на последнем наблюдении \var{KM\_Overall\_survival\_censor}, который относится к периоду после лечения и частично отражает сам исход: их включение означало бы утечку информации из будущего \parencite{kaufman2012}.

\subsection{Первичный анализ данных}

\paragraph{Проверка значений.}
Использована сводная таблица этапа извлечения признаков (298 пациентов; 20 клинических и служебных и 107 радиомических столбцов); исходный файл не изменялся, исправления вносились в копию с журналом изменений. Для 112 количественных переменных выведены наименьшее и наибольшее значения, для 14 категориальных~--- все значения в исходной записи, прочитанной без автоматического распознавания пропусков. Разных написаний одной категории нет; пропуски обозначены пустой ячейкой, а в \var{N\_category} и \var{AJCC\_Stage} встречается служебный код X (<<не может быть оценена>>). Дат в наборе нет, поэтому порядок событий и длительности не проверялись.

Правило исправления: значение заменялось верным, только если оно однозначно восстанавливается по другим полям записи; служебные коды и невосстановимые противоречия заменялись пропуском; правдоподобные, хотя и редкие значения сохранялись. Код X у одного пациента заменён пропуском в обеих переменных, так как это отметка об отсутствии данных, а не категория. У одного пациента при T4N1 указана стадия III, хотя по 7-му изданию AJCC при M0 категория T4 всегда означает стадию IV \parencite{edge2010}; стадия однозначно выводится из T и N и исправлена на IV. У 7 пациентов со статусом Current или Former указано 0 пачко-лет, что противоречит статусу; истинную экспозицию восстановить нельзя, и значения заменены пропуском. Всего изменено 10 значений: 2 в \var{AJCC\_Stage}, 1 в \var{N\_category} и 7 в \var{Smoking\_Pack-Years}.

Согласованы доза на фракцию (1,8--2,4~Гр) и отношение длительности курса к числу фракций (0,97--1,70 сут; меньше единицы при облучении дважды в день). Минимальная плотность в GTVp ниже $-1024$~HU у 9 пациентов и максимальная выше 3071~HU у 2 объясняются выбросами B-сплайн-интерполяции и артефактами от металла; это свойства вычисленных признаков, а не ошибки записи, и они сохранены. Трое пациентов с погрешностями маски, выявленными ранее, относятся к тестовой части и в анализ не входят.

\paragraph{Пропуски.}
После приведения к единой отметке пропуски есть в шести переменных: \var{Local\_tumor\_recurrence}~--- 158 (53,0~\%), \var{Pathological\_grade}~--- 43 (14,4~\%), \var{Smoking\_Pack-Years}~--- 21 (7,0~\%, из них 7 после исправления), \var{Race}, \var{N\_category} и \var{AJCC\_Stage}~--- по 1 (0,3~\%); в радиомических признаках пропусков нет. В когорте с определённым исходом остаются пропуски в \var{Pathological\_grade} (16; 11,4~\%) и \var{Smoking\_Pack-Years} (10; 7,1~\%), ни у одного пациента они не совпадают, полные данные у 114 из 140. Переменная исключалась при доле пропусков в когорте выше 30~\%, так как иначе её значения определялись бы в основном моделью заполнения, а не данными; таких переменных нет.

\paragraph{Отбор признаков.}
Исключены идентификатор \var{Patient\_ID} (сохранён как ключ строки), источник целевой переменной, служебная \var{Fold}, ставшая в когорте константой, и две переменные, недоступные в момент применения модели. Добавлена вычисленная \var{Dose\_per\_fraction} (доза, делённая на число фракций), поскольку при равной суммарной дозе биологический эффект зависит от фракционирования. Среди 123 кандидатов констант нет, а \var{GTVp\_firstorder\_TotalEnergy} совпадает с \var{GTVp\_firstorder\_Energy} во всех строках, так как при вокселе 1~мм$^3$ они отличаются лишь множителем~--- объёмом вокселя \parencite{pyradiomicsdocs}; дубликат удалён.

\paragraph{Сокращение избыточности.}
Сокращались 111 количественных признаков (5 клинических и 106 радиомических); категориальные и порядковые переменные, включая T, N и стадию, не рассматривались. Порядок приоритета не зависит от исхода: сначала клинические (\var{Age\_at\_diagnosis}, \var{Smoking\_Pack-Years}, суммарная доза, \var{Dose\_per\_fraction}, число фракций) как доступные в рутинной практике, затем радиомические по классам shape, firstorder, glcm, glrlm, glszm, gldm, ngtdm, т.\,е. от интерпретируемых к текстурным; в классе формы первыми идут \var{MeshVolume} (объём опухоли~--- установленный прогностический фактор) и \var{Sphericity}, далее порядок PyRadiomics. Признак оставлялся, если модуль коэффициента корреляции Спирмена с каждым ранее оставленным не превышал 0,9, иначе исключался и засчитывался наиболее связанному с ним оставленному. Ранговая корреляция выбрана, так как распределения признаков асимметричны, связи (например, объёма и размеров) монотонны, но нелинейны, а выбросы плотности не должны влиять на результат; порог 0,9 удаляет практически взаимозаменяемые признаки, сохраняя различные по смыслу. Корреляции считались по 140 пациентам когорты. Оставлено 43 признака (4 клинических и 39 радиомических), исключено 68, наибольший модуль корреляции среди оставленных~--- 0,8999. Больше всего признаков заменяют \var{GTVp\_firstorder\_Entropy} (14), \var{GTVp\_shape\_MeshVolume} (13) и \var{GTVp\_firstorder\_Minimum} (10); число фракций исключено как избыточное к \var{Dose\_per\_fraction}.

\paragraph{Аналитическая таблица.}
Таблица содержит 140 пациентов (строка на пациента), код пациента, целевую переменную и 54 признака: 11 категориальных, 4 количественных клинических и 39 радиомических; кодирование категорий отнесено к моделированию. На этом этапе пропуски (26 значений) не заполнялись, а связь признаков с исходом не изучалась, чтобы заполнение и отбор по исходу можно было выполнять внутри кросс-валидации без оптимистичного смещения оценки качества.

\printbibliography[title={Ссылки}]

\onecolumn
\begin{landscape}
\subsection{Словарь переменных и решений}

\noindent Словарь включает все переменные исходной таблицы, кроме радиомических, вычисленную переменную \var{Dose\_per\_fraction}, целевую переменную и 39 оставленных радиомических признаков GTVp (из 107); для каждого радиомического признака указано, сколько исключённых признаков он заменяет. В столбце «Пропуски» приведено число пропусков после исправления значений во всей таблице (298 пациентов) и в аналитической когорте (140 пациентов). Моменты доступности даны относительно момента применения модели~--- после утверждения плана лечения до его начала.

\footnotesize
\setlength{\tabcolsep}{3.5pt}
\renewcommand{\arraystretch}{1.12}
\begin{longtable}{>{\raggedright\arraybackslash}p{4.9cm} >{\raggedright\arraybackslash}p{4.2cm} >{\raggedright\arraybackslash}p{3.5cm} >{\raggedright\arraybackslash}p{3.0cm} >{\raggedright\arraybackslash}p{2.0cm} >{\raggedright\arraybackslash}p{5.9cm}}
\caption{Словарь переменных и решений первичного анализа.}\\
\toprule
\textbf{Переменная} & \textbf{Смысл} & \textbf{Тип, допустимые значения} & \textbf{Момент доступности} & \textbf{Пропуски (298 / 140)} & \textbf{Решение и комментарий} \\
\midrule
\endfirsthead
\toprule
\textbf{Переменная} & \textbf{Смысл} & \textbf{Тип, допустимые значения} & \textbf{Момент доступности} & \textbf{Пропуски (298 / 140)} & \textbf{Решение и комментарий} \\
\midrule
\endhead
\midrule
\multicolumn{6}{r}{\textit{Продолжение на следующей странице}}\\
\endfoot
\bottomrule
\endlastfoot
\texttt{Patient\_\allowbreak{}ID} & код пациента, присвоенный при анонимизации; совпадает с именем папки снимков & целый, 1–315, уникален & служебная, всегда & 0 / 0 & оставить как ключ строки; не признак: идентификатор не несёт клинической информации \\
\texttt{Fold} & часть набора соревнования & категориальный: Training, Test & служебная, задана организаторами & 0 / 0 & исключить: служебная переменная; после исключения Test — константа (Training) \\
\texttt{Local\_\allowbreak{}tumor\_\allowbreak{}recurrence} & локальный рецидив первичной опухоли за период наблюдения & бинарный: 0, 1; в Test пусто & после наблюдения не менее 2 лет & 158 (53,0~\%) / 0 & исключить из признаков: источник целевой переменной; 158 пустых значений Test — неопределённый исход \\
\texttt{Gender} & пол & категориальный: Male, Female & до начала лечения, при постановке диагноза & 0 / 0 & оставить \\
\texttt{Age\_\allowbreak{}at\_\allowbreak{}diagnosis} & возраст на момент диагноза, лет & целый, 30–89 & до начала лечения, при постановке диагноза & 0 / 0 & оставить; при сокращении избыточности не заменяет других \\
\texttt{Race} & раса & категориальный: White, Black, Hispanic, Asian, American\_Indian/\allowbreak{}Alaska\_Native & до начала лечения, при постановке диагноза & 1 (0,3~\%) / 0 & оставить; пустая ячейка — пропуск, не заполнять \\
\texttt{Tumor\_\allowbreak{}side} & латеральность опухоли & категориальный: L, R, Bilateral & до начала лечения, после стадирования & 0 / 0 & оставить \\
\texttt{Tumor\_\allowbreak{}subsite} & подобласть ротоглотки & категориальный: BOT, Tonsil, GPS, Soft\_palate, Pharyngeal\_wall, Other & до начала лечения, после стадирования & 0 / 0 & оставить \\
\texttt{T\_\allowbreak{}category} & категория T (AJCC 7) & порядковый: 1–4 & до начала лечения, после стадирования & 0 / 0 & оставить \\
\texttt{N\_\allowbreak{}category} & категория N (AJCC 7) & порядковый: 0, 1, 2a, 2b, 2c, 3 & до начала лечения, после стадирования & 1 (0,3~\%) / 0 & оставить; служебный код X (1 значение) заменён пропуском \\
\texttt{AJCC\_\allowbreak{}Stage} & стадия по AJCC 7, определяется T и N при M0 & порядковый: I–IV & до начала лечения, после стадирования & 1 (0,3~\%) / 0 & оставить; X (1) заменён пропуском, III при T4N1 (1) исправлено на IV \\
\texttt{Pathological\_\allowbreak{}grade} & степень дифференцировки опухоли & порядковый: I, I-II, II, II-III, III, IV & до начала лечения, по биопсии & 43 (14,4~\%) / 16 (11,4~\%) & оставить; пропуски не заполнять, доля ниже порога 30~\% \\
\texttt{Smoking\_\allowbreak{}status\_\allowbreak{}at\_\allowbreak{}diagnosis} & статус курения на момент диагноза & категориальный: Never, Former, Current & до начала лечения, при постановке диагноза & 0 / 0 & оставить \\
\texttt{Smoking\_\allowbreak{}Pack-Years} & кумулятивное курение, пачко-лет & непрерывный, $\geq 0$; 0 только при Never & до начала лечения, при постановке диагноза & 21 (7,0~\%) / 10 (7,1~\%) & оставить; 7 нулей при Current/Former заменены пропуском; при сокращении не заменяет других \\
\texttt{Radiation\_\allowbreak{}treatment\_\allowbreak{}course\_\allowbreak{}duration} & фактическая длительность курса ЛТ, дни & целый, 32–56 & после окончания ЛТ & 0 / 0 & исключить: неизвестна в момент применения модели \\
\texttt{Total\_\allowbreak{}prescribed\_\allowbreak{}Radiation\_\allowbreak{}treatment\_\allowbreak{}dose} & суммарная предписанная доза, Гр & целый: 60, 63, 66, 70, 72 & до начала ЛТ, при утверждении плана лечения & 0 / 0 & оставить; при сокращении не заменяет других \\
\texttt{\#\_\allowbreak{}Radiation\_\allowbreak{}treatment\_\allowbreak{}fractions} & число фракций & целый, 30–40 & до начала ЛТ, при утверждении плана лечения & 0 / 0 & исключить: избыточна, $|\rho| > 0{,}9$ с Dose\_per\_fraction \\
\texttt{Induction\_\allowbreak{}Chemotherapy} & индукционная химиотерапия до ЛТ & бинарный: Y, N & до начала ЛТ, при утверждении плана лечения & 0 / 0 & оставить \\
\texttt{Concurrent\_\allowbreak{}chemotherapy} & химиотерапия одновременно с ЛТ & бинарный: Y, N & до начала ЛТ, при утверждении плана лечения & 0 / 0 & оставить \\
\texttt{KM\_\allowbreak{}Overall\_\allowbreak{}survival\_\allowbreak{}censor} & жизненный статус на последнем наблюдении & бинарный: 1 — жив, 0 — умер & после наблюдения & 0 / 0 & исключить: неизвестна в момент применения модели, отражает исход лечения \\
\texttt{GTVn\_\allowbreak{}n\_\allowbreak{}nodes} & число поражённых лимфоузлов, выделенных на маске GTVn & целый, 0–16 & до начала лечения, по диагностической КТ & 0 / 0 & оставить; 0 означает отсутствие узлов и объясняет пропуски признаков GTVn; не заменяет других \\
\texttt{Dose\_\allowbreak{}per\_\allowbreak{}fraction} & вычисленная: доза на фракцию, Гр (доза / число фракций) & непрерывный, 1,8–2,4 & до начала ЛТ, при утверждении плана лечения & 0 / 0 & оставить; заменяет 1 исключённый: \#\_Radiation\_treatment\_fractions \\
\texttt{Target\_\allowbreak{}local\_\allowbreak{}recurrence} & целевая: 1 — локальный рецидив, 0 — нет & бинарный: 0 (128), 1 (12) & после наблюдения не менее 2 лет & 158 (53,0~\%) / 0 & целевая переменная; 158 пациентов с неопределённым исходом исключены, нулём не кодировались \\
\texttt{GTVp\_\allowbreak{}shape\_\allowbreak{}Mesh\allowbreak{}Volume} & объём по треугольной сетке (GTVp) & непрерывный, мм$^3$, $>0$ & до начала лечения, по диагностической КТ & 0 / 0 & оставить; заменяет 13 исключённых признаков: shape\_\allowbreak{}Least\allowbreak{}Axis\allowbreak{}Length, shape\_\allowbreak{}Minor\allowbreak{}Axis\allowbreak{}Length, shape\_\allowbreak{}Surface\allowbreak{}Area, shape\_\allowbreak{}Voxel\allowbreak{}Volume, firstorder\_\allowbreak{}Energy, glrlm\_\allowbreak{}Gray\allowbreak{}Level\allowbreak{}Non\allowbreak{}Uniformity, glrlm\_\allowbreak{}Run\allowbreak{}Length\allowbreak{}Non\allowbreak{}Uniformity, glszm\_\allowbreak{}Gray\allowbreak{}Level\allowbreak{}Non\allowbreak{}Uniformity, glszm\_\allowbreak{}Large\allowbreak{}Area\allowbreak{}High\allowbreak{}Gray\allowbreak{}Level\allowbreak{}Emphasis, gldm\_\allowbreak{}Dependence\allowbreak{}Non\allowbreak{}Uniformity, gldm\_\allowbreak{}Gray\allowbreak{}Level\allowbreak{}Non\allowbreak{}Uniformity, ngtdm\_\allowbreak{}Coarseness, firstorder\_\allowbreak{}Total\allowbreak{}Energy \\
\texttt{GTVp\_\allowbreak{}shape\_\allowbreak{}Sphericity} & сферичность (GTVp) & непрерывный, (0; 1] & до начала лечения, по диагностической КТ & 0 / 0 & оставить; заменяет 0 исключённых признаков \\
\texttt{GTVp\_\allowbreak{}shape\_\allowbreak{}Elongation} & вытянутость: отношение второй главной оси к первой (GTVp) & непрерывный, (0; 1] & до начала лечения, по диагностической КТ & 0 / 0 & оставить; заменяет 0 исключённых признаков \\
\texttt{GTVp\_\allowbreak{}shape\_\allowbreak{}Flatness} & уплощённость: отношение третьей главной оси к первой (GTVp) & непрерывный, (0; 1] & до начала лечения, по диагностической КТ & 0 / 0 & оставить; заменяет 0 исключённых признаков \\
\texttt{GTVp\_\allowbreak{}shape\_\allowbreak{}Major\allowbreak{}Axis\allowbreak{}Length} & длина наибольшей главной оси (GTVp) & непрерывный, мм, $>0$ & до начала лечения, по диагностической КТ & 0 / 0 & оставить; заменяет 3 исключённых признака: shape\_\allowbreak{}Maximum2DDiameter\allowbreak{}Column, shape\_\allowbreak{}Maximum2DDiameter\allowbreak{}Row, shape\_\allowbreak{}Maximum3DDiameter \\
\texttt{GTVp\_\allowbreak{}shape\_\allowbreak{}Maximum2DDiameter\allowbreak{}Slice} & наибольший диаметр в аксиальной плоскости (GTVp) & непрерывный, мм, $>0$ & до начала лечения, по диагностической КТ & 0 / 0 & оставить; заменяет 0 исключённых признаков \\
\texttt{GTVp\_\allowbreak{}shape\_\allowbreak{}Surface\allowbreak{}Volume\allowbreak{}Ratio} & отношение площади поверхности к объёму (GTVp) & непрерывный, мм$^{-1}$, $>0$ & до начала лечения, по диагностической КТ & 0 / 0 & оставить; заменяет 0 исключённых признаков \\
\texttt{GTVp\_\allowbreak{}firstorder\_\allowbreak{}10Percentile} & 10-й перцентиль плотности (GTVp) & непрерывный, HU & до начала лечения, по диагностической КТ & 0 / 0 & оставить; заменяет 0 исключённых признаков \\
\texttt{GTVp\_\allowbreak{}firstorder\_\allowbreak{}90Percentile} & 90-й перцентиль плотности (GTVp) & непрерывный, HU & до начала лечения, по диагностической КТ & 0 / 0 & оставить; заменяет 0 исключённых признаков \\
\texttt{GTVp\_\allowbreak{}firstorder\_\allowbreak{}Entropy} & энтропия гистограммы плотностей (GTVp) & непрерывный, бит, $\geq 0$ & до начала лечения, по диагностической КТ & 0 / 0 & оставить; заменяет 14 исключённых признаков: firstorder\_\allowbreak{}Interquartile\allowbreak{}Range, firstorder\_\allowbreak{}Mean\allowbreak{}Absolute\allowbreak{}Deviation, firstorder\_\allowbreak{}Robust\allowbreak{}Mean\allowbreak{}Absolute\allowbreak{}Deviation, firstorder\_\allowbreak{}Uniformity, firstorder\_\allowbreak{}Variance, glcm\_\allowbreak{}Difference\allowbreak{}Average, glcm\_\allowbreak{}Difference\allowbreak{}Entropy, glcm\_\allowbreak{}Id, glcm\_\allowbreak{}Idm, glcm\_\allowbreak{}Joint\allowbreak{}Energy, glcm\_\allowbreak{}Joint\allowbreak{}Entropy, glcm\_\allowbreak{}Sum\allowbreak{}Entropy, glrlm\_\allowbreak{}Gray\allowbreak{}Level\allowbreak{}Non\allowbreak{}Uniformity\allowbreak{}Normalized, gldm\_\allowbreak{}Dependence\allowbreak{}Entropy \\
\texttt{GTVp\_\allowbreak{}firstorder\_\allowbreak{}Kurtosis} & эксцесс распределения плотностей (GTVp) & непрерывный, $\geq 1$ & до начала лечения, по диагностической КТ & 0 / 0 & оставить; заменяет 0 исключённых признаков \\
\texttt{GTVp\_\allowbreak{}firstorder\_\allowbreak{}Maximum} & максимальная плотность (GTVp) & непрерывный, HU & до начала лечения, по диагностической КТ & 0 / 0 & оставить; заменяет 0 исключённых признаков \\
\texttt{GTVp\_\allowbreak{}firstorder\_\allowbreak{}Mean} & средняя плотность (GTVp) & непрерывный, HU & до начала лечения, по диагностической КТ & 0 / 0 & оставить; заменяет 0 исключённых признаков \\
\texttt{GTVp\_\allowbreak{}firstorder\_\allowbreak{}Median} & медиана плотности (GTVp) & непрерывный, HU & до начала лечения, по диагностической КТ & 0 / 0 & оставить; заменяет 0 исключённых признаков \\
\texttt{GTVp\_\allowbreak{}firstorder\_\allowbreak{}Minimum} & минимальная плотность (GTVp) & непрерывный, HU & до начала лечения, по диагностической КТ & 0 / 0 & оставить; заменяет 10 исключённых признаков: firstorder\_\allowbreak{}Range, glcm\_\allowbreak{}Autocorrelation, glcm\_\allowbreak{}Joint\allowbreak{}Average, glcm\_\allowbreak{}Sum\allowbreak{}Average, glrlm\_\allowbreak{}High\allowbreak{}Gray\allowbreak{}Level\allowbreak{}Run\allowbreak{}Emphasis, glrlm\_\allowbreak{}Short\allowbreak{}Run\allowbreak{}High\allowbreak{}Gray\allowbreak{}Level\allowbreak{}Emphasis, glszm\_\allowbreak{}High\allowbreak{}Gray\allowbreak{}Level\allowbreak{}Zone\allowbreak{}Emphasis, glszm\_\allowbreak{}Small\allowbreak{}Area\allowbreak{}High\allowbreak{}Gray\allowbreak{}Level\allowbreak{}Emphasis, gldm\_\allowbreak{}High\allowbreak{}Gray\allowbreak{}Level\allowbreak{}Emphasis, gldm\_\allowbreak{}Small\allowbreak{}Dependence\allowbreak{}High\allowbreak{}Gray\allowbreak{}Level\allowbreak{}Emphasis \\
\texttt{GTVp\_\allowbreak{}firstorder\_\allowbreak{}Root\allowbreak{}Mean\allowbreak{}Squared} & среднеквадратичное значение плотности (GTVp) & непрерывный, HU, $\geq 0$ & до начала лечения, по диагностической КТ & 0 / 0 & оставить; заменяет 0 исключённых признаков \\
\texttt{GTVp\_\allowbreak{}firstorder\_\allowbreak{}Skewness} & асимметрия распределения плотностей (GTVp) & непрерывный & до начала лечения, по диагностической КТ & 0 / 0 & оставить; заменяет 0 исключённых признаков \\
\texttt{GTVp\_\allowbreak{}glcm\_\allowbreak{}Cluster\allowbreak{}Prominence} & GLCM: выраженность кластеров уровней серого (GTVp) & непрерывный, $\geq 0$ & до начала лечения, по диагностической КТ & 0 / 0 & оставить; заменяет 9 исключённых признаков: glcm\_\allowbreak{}Cluster\allowbreak{}Tendency, glcm\_\allowbreak{}Contrast, glcm\_\allowbreak{}Difference\allowbreak{}Variance, glcm\_\allowbreak{}Sum\allowbreak{}Squares, glrlm\_\allowbreak{}Gray\allowbreak{}Level\allowbreak{}Variance, glszm\_\allowbreak{}Gray\allowbreak{}Level\allowbreak{}Non\allowbreak{}Uniformity\allowbreak{}Normalized, glszm\_\allowbreak{}Gray\allowbreak{}Level\allowbreak{}Variance, gldm\_\allowbreak{}Gray\allowbreak{}Level\allowbreak{}Variance, ngtdm\_\allowbreak{}Complexity \\
\texttt{GTVp\_\allowbreak{}glcm\_\allowbreak{}Cluster\allowbreak{}Shade} & GLCM: асимметрия кластеров (GTVp) & непрерывный & до начала лечения, по диагностической КТ & 0 / 0 & оставить; заменяет 0 исключённых признаков \\
\texttt{GTVp\_\allowbreak{}glcm\_\allowbreak{}Correlation} & GLCM: корреляция уровней соседних вокселей (GTVp) & непрерывный, $[-1; 1]$ & до начала лечения, по диагностической КТ & 0 / 0 & оставить; заменяет 0 исключённых признаков \\
\texttt{GTVp\_\allowbreak{}glcm\_\allowbreak{}Idmn} & GLCM: нормированный обратный момент разности (однородность) (GTVp) & непрерывный, (0; 1] & до начала лечения, по диагностической КТ & 0 / 0 & оставить; заменяет 1 исключённый признак: glcm\_\allowbreak{}Idn \\
\texttt{GTVp\_\allowbreak{}glcm\_\allowbreak{}Imc1} & GLCM: информационная мера корреляции 1 (GTVp) & непрерывный, $[-1; 0]$ & до начала лечения, по диагностической КТ & 0 / 0 & оставить; заменяет 1 исключённый признак: glcm\_\allowbreak{}Imc2 \\
\texttt{GTVp\_\allowbreak{}glcm\_\allowbreak{}Inverse\allowbreak{}Variance} & GLCM: обратная дисперсия (GTVp) & непрерывный, $\geq 0$ & до начала лечения, по диагностической КТ & 0 / 0 & оставить; заменяет 0 исключённых признаков \\
\texttt{GTVp\_\allowbreak{}glcm\_\allowbreak{}MCC} & GLCM: максимальный коэффициент корреляции (GTVp) & непрерывный, [0; 1] & до начала лечения, по диагностической КТ & 0 / 0 & оставить; заменяет 0 исключённых признаков \\
\texttt{GTVp\_\allowbreak{}glcm\_\allowbreak{}Maximum\allowbreak{}Probability} & GLCM: частота наиболее частой пары уровней (GTVp) & непрерывный, (0; 1] & до начала лечения, по диагностической КТ & 0 / 0 & оставить; заменяет 0 исключённых признаков \\
\texttt{GTVp\_\allowbreak{}glrlm\_\allowbreak{}Long\allowbreak{}Run\allowbreak{}Emphasis} & GLRLM: преобладание длинных серий (GTVp) & непрерывный, $\geq 1$ & до начала лечения, по диагностической КТ & 0 / 0 & оставить; заменяет 5 исключённых признаков: glrlm\_\allowbreak{}Run\allowbreak{}Length\allowbreak{}Non\allowbreak{}Uniformity\allowbreak{}Normalized, glrlm\_\allowbreak{}Run\allowbreak{}Percentage, glrlm\_\allowbreak{}Run\allowbreak{}Variance, glrlm\_\allowbreak{}Short\allowbreak{}Run\allowbreak{}Emphasis, gldm\_\allowbreak{}Large\allowbreak{}Dependence\allowbreak{}Emphasis \\
\texttt{GTVp\_\allowbreak{}glrlm\_\allowbreak{}Long\allowbreak{}Run\allowbreak{}High\allowbreak{}Gray\allowbreak{}Level\allowbreak{}Emphasis} & GLRLM: длинные серии высоких уровней (GTVp) & непрерывный, $>0$ & до начала лечения, по диагностической КТ & 0 / 0 & оставить; заменяет 1 исключённый признак: gldm\_\allowbreak{}Large\allowbreak{}Dependence\allowbreak{}High\allowbreak{}Gray\allowbreak{}Level\allowbreak{}Emphasis \\
\texttt{GTVp\_\allowbreak{}glrlm\_\allowbreak{}Long\allowbreak{}Run\allowbreak{}Low\allowbreak{}Gray\allowbreak{}Level\allowbreak{}Emphasis} & GLRLM: длинные серии низких уровней (GTVp) & непрерывный, $>0$ & до начала лечения, по диагностической КТ & 0 / 0 & оставить; заменяет 1 исключённый признак: gldm\_\allowbreak{}Large\allowbreak{}Dependence\allowbreak{}Low\allowbreak{}Gray\allowbreak{}Level\allowbreak{}Emphasis \\
\texttt{GTVp\_\allowbreak{}glrlm\_\allowbreak{}Low\allowbreak{}Gray\allowbreak{}Level\allowbreak{}Run\allowbreak{}Emphasis} & GLRLM: преобладание серий низких уровней (GTVp) & непрерывный, (0; 1] & до начала лечения, по диагностической КТ & 0 / 0 & оставить; заменяет 3 исключённых признака: glrlm\_\allowbreak{}Short\allowbreak{}Run\allowbreak{}Low\allowbreak{}Gray\allowbreak{}Level\allowbreak{}Emphasis, glszm\_\allowbreak{}Low\allowbreak{}Gray\allowbreak{}Level\allowbreak{}Zone\allowbreak{}Emphasis, gldm\_\allowbreak{}Low\allowbreak{}Gray\allowbreak{}Level\allowbreak{}Emphasis \\
\texttt{GTVp\_\allowbreak{}glrlm\_\allowbreak{}Run\allowbreak{}Entropy} & GLRLM: энтропия серий (GTVp) & непрерывный, бит, $\geq 0$ & до начала лечения, по диагностической КТ & 0 / 0 & оставить; заменяет 1 исключённый признак: glszm\_\allowbreak{}Zone\allowbreak{}Entropy \\
\texttt{GTVp\_\allowbreak{}glszm\_\allowbreak{}Large\allowbreak{}Area\allowbreak{}Emphasis} & GLSZM: преобладание крупных зон (GTVp) & непрерывный, $\geq 1$ & до начала лечения, по диагностической КТ & 0 / 0 & оставить; заменяет 1 исключённый признак: glszm\_\allowbreak{}Zone\allowbreak{}Variance \\
\texttt{GTVp\_\allowbreak{}glszm\_\allowbreak{}Large\allowbreak{}Area\allowbreak{}Low\allowbreak{}Gray\allowbreak{}Level\allowbreak{}Emphasis} & GLSZM: крупные зоны низких уровней (GTVp) & непрерывный, $>0$ & до начала лечения, по диагностической КТ & 0 / 0 & оставить; заменяет 2 исключённых признака: glszm\_\allowbreak{}Zone\allowbreak{}Percentage, ngtdm\_\allowbreak{}Busyness \\
\texttt{GTVp\_\allowbreak{}glszm\_\allowbreak{}Size\allowbreak{}Zone\allowbreak{}Non\allowbreak{}Uniformity} & GLSZM: неоднородность размеров зон (GTVp) & непрерывный, $>0$ & до начала лечения, по диагностической КТ & 0 / 0 & оставить; заменяет 1 исключённый признак: glszm\_\allowbreak{}Small\allowbreak{}Area\allowbreak{}Low\allowbreak{}Gray\allowbreak{}Level\allowbreak{}Emphasis \\
\texttt{GTVp\_\allowbreak{}glszm\_\allowbreak{}Size\allowbreak{}Zone\allowbreak{}Non\allowbreak{}Uniformity\allowbreak{}Normalized} & GLSZM: нормированная неоднородность размеров зон (GTVp) & непрерывный, (0; 1] & до начала лечения, по диагностической КТ & 0 / 0 & оставить; заменяет 1 исключённый признак: glszm\_\allowbreak{}Small\allowbreak{}Area\allowbreak{}Emphasis \\
\texttt{GTVp\_\allowbreak{}gldm\_\allowbreak{}Dependence\allowbreak{}Non\allowbreak{}Uniformity\allowbreak{}Normalized} & GLDM: нормированная неоднородность зависимостей (GTVp) & непрерывный, (0; 1] & до начала лечения, по диагностической КТ & 0 / 0 & оставить; заменяет 1 исключённый признак: gldm\_\allowbreak{}Dependence\allowbreak{}Variance \\
\texttt{GTVp\_\allowbreak{}gldm\_\allowbreak{}Small\allowbreak{}Dependence\allowbreak{}Emphasis} & GLDM: преобладание малых зависимостей (GTVp) & непрерывный, (0; 1] & до начала лечения, по диагностической КТ & 0 / 0 & оставить; заменяет 0 исключённых признаков \\
\texttt{GTVp\_\allowbreak{}gldm\_\allowbreak{}Small\allowbreak{}Dependence\allowbreak{}Low\allowbreak{}Gray\allowbreak{}Level\allowbreak{}Emphasis} & GLDM: малые зависимости низких уровней (GTVp) & непрерывный, (0; 1] & до начала лечения, по диагностической КТ & 0 / 0 & оставить; заменяет 0 исключённых признаков \\
\texttt{GTVp\_\allowbreak{}ngtdm\_\allowbreak{}Contrast} & NGTDM: контраст соседних тонов (GTVp) & непрерывный, $\geq 0$ & до начала лечения, по диагностической КТ & 0 / 0 & оставить; заменяет 0 исключённых признаков \\
\texttt{GTVp\_\allowbreak{}ngtdm\_\allowbreak{}Strength} & NGTDM: сила текстуры (GTVp) & непрерывный, $\geq 0$ & до начала лечения, по диагностической КТ & 0 / 0 & оставить; заменяет 0 исключённых признаков \\
\texttt{GTVn\_\allowbreak{}shape\_\allowbreak{}Mesh\allowbreak{}Volume} & объём по треугольной сетке (GTVn) & непрерывный, мм$^3$, $>0$ & до начала лечения, по диагностической КТ & 45 (15,1~\%) / 20 (14,3~\%) & оставить; заменяет 7 исключённых признаков: shape\_\allowbreak{}Surface\allowbreak{}Area, shape\_\allowbreak{}Voxel\allowbreak{}Volume, glrlm\_\allowbreak{}Gray\allowbreak{}Level\allowbreak{}Non\allowbreak{}Uniformity, glrlm\_\allowbreak{}Run\allowbreak{}Length\allowbreak{}Non\allowbreak{}Uniformity, gldm\_\allowbreak{}Dependence\allowbreak{}Non\allowbreak{}Uniformity, gldm\_\allowbreak{}Gray\allowbreak{}Level\allowbreak{}Non\allowbreak{}Uniformity, ngtdm\_\allowbreak{}Coarseness; пропуск — нет узлов (GTVn\_n\_nodes = 0), не заполнять \\
\texttt{GTVn\_\allowbreak{}shape\_\allowbreak{}Sphericity} & сферичность (GTVn) & непрерывный, (0; 1] & до начала лечения, по диагностической КТ & 45 (15,1~\%) / 20 (14,3~\%) & оставить; заменяет 0 исключённых признаков; пропуск — нет узлов (GTVn\_n\_nodes = 0), не заполнять \\
\texttt{GTVn\_\allowbreak{}shape\_\allowbreak{}Elongation} & вытянутость: отношение второй главной оси к первой (GTVn) & непрерывный, (0; 1] & до начала лечения, по диагностической КТ & 45 (15,1~\%) / 20 (14,3~\%) & оставить; заменяет 0 исключённых признаков; пропуск — нет узлов (GTVn\_n\_nodes = 0), не заполнять \\
\texttt{GTVn\_\allowbreak{}shape\_\allowbreak{}Flatness} & уплощённость: отношение третьей главной оси к первой (GTVn) & непрерывный, (0; 1] & до начала лечения, по диагностической КТ & 45 (15,1~\%) / 20 (14,3~\%) & оставить; заменяет 0 исключённых признаков; пропуск — нет узлов (GTVn\_n\_nodes = 0), не заполнять \\
\texttt{GTVn\_\allowbreak{}shape\_\allowbreak{}Least\allowbreak{}Axis\allowbreak{}Length} & длина наименьшей главной оси (GTVn) & непрерывный, мм, $>0$ & до начала лечения, по диагностической КТ & 45 (15,1~\%) / 20 (14,3~\%) & оставить; заменяет 0 исключённых признаков; пропуск — нет узлов (GTVn\_n\_nodes = 0), не заполнять \\
\texttt{GTVn\_\allowbreak{}shape\_\allowbreak{}Major\allowbreak{}Axis\allowbreak{}Length} & длина наибольшей главной оси (GTVn) & непрерывный, мм, $>0$ & до начала лечения, по диагностической КТ & 45 (15,1~\%) / 20 (14,3~\%) & оставить; заменяет 2 исключённых признака: shape\_\allowbreak{}Maximum2DDiameter\allowbreak{}Column, shape\_\allowbreak{}Maximum3DDiameter; пропуск — нет узлов (GTVn\_n\_nodes = 0), не заполнять \\
\texttt{GTVn\_\allowbreak{}shape\_\allowbreak{}Maximum2DDiameter\allowbreak{}Row} & наибольший диаметр в сагиттальной плоскости (GTVn) & непрерывный, мм, $>0$ & до начала лечения, по диагностической КТ & 45 (15,1~\%) / 20 (14,3~\%) & оставить; заменяет 0 исключённых признаков; пропуск — нет узлов (GTVn\_n\_nodes = 0), не заполнять \\
\texttt{GTVn\_\allowbreak{}shape\_\allowbreak{}Maximum2DDiameter\allowbreak{}Slice} & наибольший диаметр в аксиальной плоскости (GTVn) & непрерывный, мм, $>0$ & до начала лечения, по диагностической КТ & 45 (15,1~\%) / 20 (14,3~\%) & оставить; заменяет 1 исключённый признак: shape\_\allowbreak{}Minor\allowbreak{}Axis\allowbreak{}Length; пропуск — нет узлов (GTVn\_n\_nodes = 0), не заполнять \\
\texttt{GTVn\_\allowbreak{}shape\_\allowbreak{}Surface\allowbreak{}Volume\allowbreak{}Ratio} & отношение площади поверхности к объёму (GTVn) & непрерывный, мм$^{-1}$, $>0$ & до начала лечения, по диагностической КТ & 45 (15,1~\%) / 20 (14,3~\%) & оставить; заменяет 1 исключённый признак: glszm\_\allowbreak{}Large\allowbreak{}Area\allowbreak{}Emphasis; пропуск — нет узлов (GTVn\_n\_nodes = 0), не заполнять \\
\texttt{GTVn\_\allowbreak{}firstorder\_\allowbreak{}10Percentile} & 10-й перцентиль плотности (GTVn) & непрерывный, HU & до начала лечения, по диагностической КТ & 45 (15,1~\%) / 20 (14,3~\%) & оставить; заменяет 0 исключённых признаков; пропуск — нет узлов (GTVn\_n\_nodes = 0), не заполнять \\
\texttt{GTVn\_\allowbreak{}firstorder\_\allowbreak{}90Percentile} & 90-й перцентиль плотности (GTVn) & непрерывный, HU & до начала лечения, по диагностической КТ & 45 (15,1~\%) / 20 (14,3~\%) & оставить; заменяет 0 исключённых признаков; пропуск — нет узлов (GTVn\_n\_nodes = 0), не заполнять \\
\texttt{GTVn\_\allowbreak{}firstorder\_\allowbreak{}Energy} & энергия: сумма квадратов плотностей (GTVn) & непрерывный, $\geq 0$ & до начала лечения, по диагностической КТ & 45 (15,1~\%) / 20 (14,3~\%) & оставить; заменяет 1 исключённый признак: firstorder\_\allowbreak{}Total\allowbreak{}Energy; пропуск — нет узлов (GTVn\_n\_nodes = 0), не заполнять \\
\texttt{GTVn\_\allowbreak{}firstorder\_\allowbreak{}Entropy} & энтропия гистограммы плотностей (GTVn) & непрерывный, бит, $\geq 0$ & до начала лечения, по диагностической КТ & 45 (15,1~\%) / 20 (14,3~\%) & оставить; заменяет 12 исключённых признаков: firstorder\_\allowbreak{}Interquartile\allowbreak{}Range, firstorder\_\allowbreak{}Mean\allowbreak{}Absolute\allowbreak{}Deviation, firstorder\_\allowbreak{}Robust\allowbreak{}Mean\allowbreak{}Absolute\allowbreak{}Deviation, firstorder\_\allowbreak{}Uniformity, firstorder\_\allowbreak{}Variance, glcm\_\allowbreak{}Cluster\allowbreak{}Tendency, glcm\_\allowbreak{}Joint\allowbreak{}Energy, glcm\_\allowbreak{}Joint\allowbreak{}Entropy, glcm\_\allowbreak{}Sum\allowbreak{}Entropy, glcm\_\allowbreak{}Sum\allowbreak{}Squares, glrlm\_\allowbreak{}Gray\allowbreak{}Level\allowbreak{}Non\allowbreak{}Uniformity\allowbreak{}Normalized, gldm\_\allowbreak{}Gray\allowbreak{}Level\allowbreak{}Variance; пропуск — нет узлов (GTVn\_n\_nodes = 0), не заполнять \\
\texttt{GTVn\_\allowbreak{}firstorder\_\allowbreak{}Kurtosis} & эксцесс распределения плотностей (GTVn) & непрерывный, $\geq 1$ & до начала лечения, по диагностической КТ & 45 (15,1~\%) / 20 (14,3~\%) & оставить; заменяет 0 исключённых признаков; пропуск — нет узлов (GTVn\_n\_nodes = 0), не заполнять \\
\texttt{GTVn\_\allowbreak{}firstorder\_\allowbreak{}Maximum} & максимальная плотность (GTVn) & непрерывный, HU & до начала лечения, по диагностической КТ & 45 (15,1~\%) / 20 (14,3~\%) & оставить; заменяет 0 исключённых признаков; пропуск — нет узлов (GTVn\_n\_nodes = 0), не заполнять \\
\texttt{GTVn\_\allowbreak{}firstorder\_\allowbreak{}Mean} & средняя плотность (GTVn) & непрерывный, HU & до начала лечения, по диагностической КТ & 45 (15,1~\%) / 20 (14,3~\%) & оставить; заменяет 2 исключённых признака: firstorder\_\allowbreak{}Median, firstorder\_\allowbreak{}Root\allowbreak{}Mean\allowbreak{}Squared; пропуск — нет узлов (GTVn\_n\_nodes = 0), не заполнять \\
\texttt{GTVn\_\allowbreak{}firstorder\_\allowbreak{}Minimum} & минимальная плотность (GTVn) & непрерывный, HU & до начала лечения, по диагностической КТ & 45 (15,1~\%) / 20 (14,3~\%) & оставить; заменяет 0 исключённых признаков; пропуск — нет узлов (GTVn\_n\_nodes = 0), не заполнять \\
\texttt{GTVn\_\allowbreak{}firstorder\_\allowbreak{}Range} & размах плотностей (GTVn) & непрерывный, HU, $\geq 0$ & до начала лечения, по диагностической КТ & 45 (15,1~\%) / 20 (14,3~\%) & оставить; заменяет 1 исключённый признак: ngtdm\_\allowbreak{}Complexity; пропуск — нет узлов (GTVn\_n\_nodes = 0), не заполнять \\
\texttt{GTVn\_\allowbreak{}firstorder\_\allowbreak{}Skewness} & асимметрия распределения плотностей (GTVn) & непрерывный & до начала лечения, по диагностической КТ & 45 (15,1~\%) / 20 (14,3~\%) & оставить; заменяет 0 исключённых признаков; пропуск — нет узлов (GTVn\_n\_nodes = 0), не заполнять \\
\texttt{GTVn\_\allowbreak{}glcm\_\allowbreak{}Autocorrelation} & GLCM: автокорреляция уровней серого (GTVn) & непрерывный, $>0$ & до начала лечения, по диагностической КТ & 45 (15,1~\%) / 20 (14,3~\%) & оставить; заменяет 5 исключённых признаков: glcm\_\allowbreak{}Joint\allowbreak{}Average, glcm\_\allowbreak{}Sum\allowbreak{}Average, glrlm\_\allowbreak{}High\allowbreak{}Gray\allowbreak{}Level\allowbreak{}Run\allowbreak{}Emphasis, glrlm\_\allowbreak{}Short\allowbreak{}Run\allowbreak{}High\allowbreak{}Gray\allowbreak{}Level\allowbreak{}Emphasis, gldm\_\allowbreak{}High\allowbreak{}Gray\allowbreak{}Level\allowbreak{}Emphasis; пропуск — нет узлов (GTVn\_n\_nodes = 0), не заполнять \\
\texttt{GTVn\_\allowbreak{}glcm\_\allowbreak{}Cluster\allowbreak{}Prominence} & GLCM: выраженность кластеров уровней серого (GTVn) & непрерывный, $\geq 0$ & до начала лечения, по диагностической КТ & 45 (15,1~\%) / 20 (14,3~\%) & оставить; заменяет 1 исключённый признак: glrlm\_\allowbreak{}Gray\allowbreak{}Level\allowbreak{}Variance; пропуск — нет узлов (GTVn\_n\_nodes = 0), не заполнять \\
\texttt{GTVn\_\allowbreak{}glcm\_\allowbreak{}Cluster\allowbreak{}Shade} & GLCM: асимметрия кластеров (GTVn) & непрерывный & до начала лечения, по диагностической КТ & 45 (15,1~\%) / 20 (14,3~\%) & оставить; заменяет 0 исключённых признаков; пропуск — нет узлов (GTVn\_n\_nodes = 0), не заполнять \\
\texttt{GTVn\_\allowbreak{}glcm\_\allowbreak{}Contrast} & GLCM: контраст соседних вокселей (GTVn) & непрерывный, $\geq 0$ & до начала лечения, по диагностической КТ & 45 (15,1~\%) / 20 (14,3~\%) & оставить; заменяет 9 исключённых признаков: glcm\_\allowbreak{}Difference\allowbreak{}Average, glcm\_\allowbreak{}Difference\allowbreak{}Entropy, glcm\_\allowbreak{}Difference\allowbreak{}Variance, glcm\_\allowbreak{}Id, glcm\_\allowbreak{}Idm, glrlm\_\allowbreak{}Long\allowbreak{}Run\allowbreak{}Emphasis, glrlm\_\allowbreak{}Run\allowbreak{}Length\allowbreak{}Non\allowbreak{}Uniformity\allowbreak{}Normalized, glrlm\_\allowbreak{}Run\allowbreak{}Percentage, glrlm\_\allowbreak{}Short\allowbreak{}Run\allowbreak{}Emphasis; пропуск — нет узлов (GTVn\_n\_nodes = 0), не заполнять \\
\texttt{GTVn\_\allowbreak{}glcm\_\allowbreak{}Correlation} & GLCM: корреляция уровней соседних вокселей (GTVn) & непрерывный, $[-1; 1]$ & до начала лечения, по диагностической КТ & 45 (15,1~\%) / 20 (14,3~\%) & оставить; заменяет 1 исключённый признак: glcm\_\allowbreak{}Imc2; пропуск — нет узлов (GTVn\_n\_nodes = 0), не заполнять \\
\texttt{GTVn\_\allowbreak{}glcm\_\allowbreak{}Idmn} & GLCM: нормированный обратный момент разности (однородность) (GTVn) & непрерывный, (0; 1] & до начала лечения, по диагностической КТ & 45 (15,1~\%) / 20 (14,3~\%) & оставить; заменяет 1 исключённый признак: glcm\_\allowbreak{}Idn; пропуск — нет узлов (GTVn\_n\_nodes = 0), не заполнять \\
\texttt{GTVn\_\allowbreak{}glcm\_\allowbreak{}Imc1} & GLCM: информационная мера корреляции 1 (GTVn) & непрерывный, $[-1; 0]$ & до начала лечения, по диагностической КТ & 45 (15,1~\%) / 20 (14,3~\%) & оставить; заменяет 0 исключённых признаков; пропуск — нет узлов (GTVn\_n\_nodes = 0), не заполнять \\
\texttt{GTVn\_\allowbreak{}glcm\_\allowbreak{}Inverse\allowbreak{}Variance} & GLCM: обратная дисперсия (GTVn) & непрерывный, $\geq 0$ & до начала лечения, по диагностической КТ & 45 (15,1~\%) / 20 (14,3~\%) & оставить; заменяет 1 исключённый признак: gldm\_\allowbreak{}Dependence\allowbreak{}Variance; пропуск — нет узлов (GTVn\_n\_nodes = 0), не заполнять \\
\texttt{GTVn\_\allowbreak{}glcm\_\allowbreak{}MCC} & GLCM: максимальный коэффициент корреляции (GTVn) & непрерывный, [0; 1] & до начала лечения, по диагностической КТ & 45 (15,1~\%) / 20 (14,3~\%) & оставить; заменяет 0 исключённых признаков; пропуск — нет узлов (GTVn\_n\_nodes = 0), не заполнять \\
\texttt{GTVn\_\allowbreak{}glcm\_\allowbreak{}Maximum\allowbreak{}Probability} & GLCM: частота наиболее частой пары уровней (GTVn) & непрерывный, (0; 1] & до начала лечения, по диагностической КТ & 45 (15,1~\%) / 20 (14,3~\%) & оставить; заменяет 0 исключённых признаков; пропуск — нет узлов (GTVn\_n\_nodes = 0), не заполнять \\
\texttt{GTVn\_\allowbreak{}glrlm\_\allowbreak{}Long\allowbreak{}Run\allowbreak{}High\allowbreak{}Gray\allowbreak{}Level\allowbreak{}Emphasis} & GLRLM: длинные серии высоких уровней (GTVn) & непрерывный, $>0$ & до начала лечения, по диагностической КТ & 45 (15,1~\%) / 20 (14,3~\%) & оставить; заменяет 1 исключённый признак: gldm\_\allowbreak{}Large\allowbreak{}Dependence\allowbreak{}High\allowbreak{}Gray\allowbreak{}Level\allowbreak{}Emphasis; пропуск — нет узлов (GTVn\_n\_nodes = 0), не заполнять \\
\texttt{GTVn\_\allowbreak{}glrlm\_\allowbreak{}Long\allowbreak{}Run\allowbreak{}Low\allowbreak{}Gray\allowbreak{}Level\allowbreak{}Emphasis} & GLRLM: длинные серии низких уровней (GTVn) & непрерывный, $>0$ & до начала лечения, по диагностической КТ & 45 (15,1~\%) / 20 (14,3~\%) & оставить; заменяет 2 исключённых признака: gldm\_\allowbreak{}Large\allowbreak{}Dependence\allowbreak{}Low\allowbreak{}Gray\allowbreak{}Level\allowbreak{}Emphasis, gldm\_\allowbreak{}Small\allowbreak{}Dependence\allowbreak{}High\allowbreak{}Gray\allowbreak{}Level\allowbreak{}Emphasis; пропуск — нет узлов (GTVn\_n\_nodes = 0), не заполнять \\
\texttt{GTVn\_\allowbreak{}glrlm\_\allowbreak{}Low\allowbreak{}Gray\allowbreak{}Level\allowbreak{}Run\allowbreak{}Emphasis} & GLRLM: преобладание серий низких уровней (GTVn) & непрерывный, (0; 1] & до начала лечения, по диагностической КТ & 45 (15,1~\%) / 20 (14,3~\%) & оставить; заменяет 3 исключённых признака: glrlm\_\allowbreak{}Short\allowbreak{}Run\allowbreak{}Low\allowbreak{}Gray\allowbreak{}Level\allowbreak{}Emphasis, glszm\_\allowbreak{}Low\allowbreak{}Gray\allowbreak{}Level\allowbreak{}Zone\allowbreak{}Emphasis, gldm\_\allowbreak{}Low\allowbreak{}Gray\allowbreak{}Level\allowbreak{}Emphasis; пропуск — нет узлов (GTVn\_n\_nodes = 0), не заполнять \\
\texttt{GTVn\_\allowbreak{}glrlm\_\allowbreak{}Run\allowbreak{}Entropy} & GLRLM: энтропия серий (GTVn) & непрерывный, бит, $\geq 0$ & до начала лечения, по диагностической КТ & 45 (15,1~\%) / 20 (14,3~\%) & оставить; заменяет 0 исключённых признаков; пропуск — нет узлов (GTVn\_n\_nodes = 0), не заполнять \\
\texttt{GTVn\_\allowbreak{}glrlm\_\allowbreak{}Run\allowbreak{}Variance} & GLRLM: дисперсия длин серий (GTVn) & непрерывный, $\geq 0$ & до начала лечения, по диагностической КТ & 45 (15,1~\%) / 20 (14,3~\%) & оставить; заменяет 1 исключённый признак: gldm\_\allowbreak{}Large\allowbreak{}Dependence\allowbreak{}Emphasis; пропуск — нет узлов (GTVn\_n\_nodes = 0), не заполнять \\
\texttt{GTVn\_\allowbreak{}glszm\_\allowbreak{}Gray\allowbreak{}Level\allowbreak{}Non\allowbreak{}Uniformity} & GLSZM: неоднородность уровней серого по зонам (GTVn) & непрерывный, $>0$ & до начала лечения, по диагностической КТ & 45 (15,1~\%) / 20 (14,3~\%) & оставить; заменяет 1 исключённый признак: glszm\_\allowbreak{}Size\allowbreak{}Zone\allowbreak{}Non\allowbreak{}Uniformity; пропуск — нет узлов (GTVn\_n\_nodes = 0), не заполнять \\
\texttt{GTVn\_\allowbreak{}glszm\_\allowbreak{}Gray\allowbreak{}Level\allowbreak{}Non\allowbreak{}Uniformity\allowbreak{}Normalized} & GLSZM: нормированная неоднородность уровней (GTVn) & непрерывный, (0; 1] & до начала лечения, по диагностической КТ & 45 (15,1~\%) / 20 (14,3~\%) & оставить; заменяет 1 исключённый признак: glszm\_\allowbreak{}Gray\allowbreak{}Level\allowbreak{}Variance; пропуск — нет узлов (GTVn\_n\_nodes = 0), не заполнять \\
\texttt{GTVn\_\allowbreak{}glszm\_\allowbreak{}High\allowbreak{}Gray\allowbreak{}Level\allowbreak{}Zone\allowbreak{}Emphasis} & GLSZM: преобладание зон высоких уровней (GTVn) & непрерывный, $\geq 1$ & до начала лечения, по диагностической КТ & 45 (15,1~\%) / 20 (14,3~\%) & оставить; заменяет 1 исключённый признак: glszm\_\allowbreak{}Small\allowbreak{}Area\allowbreak{}High\allowbreak{}Gray\allowbreak{}Level\allowbreak{}Emphasis; пропуск — нет узлов (GTVn\_n\_nodes = 0), не заполнять \\
\texttt{GTVn\_\allowbreak{}glszm\_\allowbreak{}Large\allowbreak{}Area\allowbreak{}High\allowbreak{}Gray\allowbreak{}Level\allowbreak{}Emphasis} & GLSZM: крупные зоны высоких уровней (GTVn) & непрерывный, $>0$ & до начала лечения, по диагностической КТ & 45 (15,1~\%) / 20 (14,3~\%) & оставить; заменяет 1 исключённый признак: glszm\_\allowbreak{}Zone\allowbreak{}Variance; пропуск — нет узлов (GTVn\_n\_nodes = 0), не заполнять \\
\texttt{GTVn\_\allowbreak{}glszm\_\allowbreak{}Large\allowbreak{}Area\allowbreak{}Low\allowbreak{}Gray\allowbreak{}Level\allowbreak{}Emphasis} & GLSZM: крупные зоны низких уровней (GTVn) & непрерывный, $>0$ & до начала лечения, по диагностической КТ & 45 (15,1~\%) / 20 (14,3~\%) & оставить; заменяет 2 исключённых признака: glszm\_\allowbreak{}Zone\allowbreak{}Percentage, gldm\_\allowbreak{}Small\allowbreak{}Dependence\allowbreak{}Emphasis; пропуск — нет узлов (GTVn\_n\_nodes = 0), не заполнять \\
\texttt{GTVn\_\allowbreak{}glszm\_\allowbreak{}Size\allowbreak{}Zone\allowbreak{}Non\allowbreak{}Uniformity\allowbreak{}Normalized} & GLSZM: нормированная неоднородность размеров зон (GTVn) & непрерывный, (0; 1] & до начала лечения, по диагностической КТ & 45 (15,1~\%) / 20 (14,3~\%) & оставить; заменяет 1 исключённый признак: glszm\_\allowbreak{}Small\allowbreak{}Area\allowbreak{}Emphasis; пропуск — нет узлов (GTVn\_n\_nodes = 0), не заполнять \\
\texttt{GTVn\_\allowbreak{}glszm\_\allowbreak{}Small\allowbreak{}Area\allowbreak{}Low\allowbreak{}Gray\allowbreak{}Level\allowbreak{}Emphasis} & GLSZM: мелкие зоны низких уровней (GTVn) & непрерывный, (0; 1] & до начала лечения, по диагностической КТ & 45 (15,1~\%) / 20 (14,3~\%) & оставить; заменяет 0 исключённых признаков; пропуск — нет узлов (GTVn\_n\_nodes = 0), не заполнять \\
\texttt{GTVn\_\allowbreak{}glszm\_\allowbreak{}Zone\allowbreak{}Entropy} & GLSZM: энтропия зон (GTVn) & непрерывный, бит, $\geq 0$ & до начала лечения, по диагностической КТ & 45 (15,1~\%) / 20 (14,3~\%) & оставить; заменяет 0 исключённых признаков; пропуск — нет узлов (GTVn\_n\_nodes = 0), не заполнять \\
\texttt{GTVn\_\allowbreak{}gldm\_\allowbreak{}Dependence\allowbreak{}Entropy} & GLDM: энтропия зависимостей (GTVn) & непрерывный, бит, $\geq 0$ & до начала лечения, по диагностической КТ & 45 (15,1~\%) / 20 (14,3~\%) & оставить; заменяет 0 исключённых признаков; пропуск — нет узлов (GTVn\_n\_nodes = 0), не заполнять \\
\texttt{GTVn\_\allowbreak{}gldm\_\allowbreak{}Dependence\allowbreak{}Non\allowbreak{}Uniformity\allowbreak{}Normalized} & GLDM: нормированная неоднородность зависимостей (GTVn) & непрерывный, (0; 1] & до начала лечения, по диагностической КТ & 45 (15,1~\%) / 20 (14,3~\%) & оставить; заменяет 0 исключённых признаков; пропуск — нет узлов (GTVn\_n\_nodes = 0), не заполнять \\
\texttt{GTVn\_\allowbreak{}gldm\_\allowbreak{}Small\allowbreak{}Dependence\allowbreak{}Low\allowbreak{}Gray\allowbreak{}Level\allowbreak{}Emphasis} & GLDM: малые зависимости низких уровней (GTVn) & непрерывный, (0; 1] & до начала лечения, по диагностической КТ & 45 (15,1~\%) / 20 (14,3~\%) & оставить; заменяет 0 исключённых признаков; пропуск — нет узлов (GTVn\_n\_nodes = 0), не заполнять \\
\texttt{GTVn\_\allowbreak{}ngtdm\_\allowbreak{}Busyness} & NGTDM: «занятость» — частота смены тонов (GTVn) & непрерывный, $\geq 0$ & до начала лечения, по диагностической КТ & 45 (15,1~\%) / 20 (14,3~\%) & оставить; заменяет 1 исключённый признак: ngtdm\_\allowbreak{}Strength; пропуск — нет узлов (GTVn\_n\_nodes = 0), не заполнять \\
\texttt{GTVn\_\allowbreak{}ngtdm\_\allowbreak{}Contrast} & NGTDM: контраст соседних тонов (GTVn) & непрерывный, $\geq 0$ & до начала лечения, по диагностической КТ & 45 (15,1~\%) / 20 (14,3~\%) & оставить; заменяет 0 исключённых признаков; пропуск — нет узлов (GTVn\_n\_nodes = 0), не заполнять \\

\end{longtable}
\end{landscape}

\end{document}
